\subsubsection{Symbian OS}

\sourcenotebox{\textbf{Source Note:} The quoted listings below are taken from
the local Symbian EPOC Kernel Architecture 2 (EKA2) source tree in this
workspace under \texttt{kernel/eka/}.
The upstream repository is
\url{https://github.com/SymbianSource/oss.FCL.sf.os.kernelhwsrv}. Each caption
names the relevant source file or files.}

Symbian EKA2 uses a \textbf{fixed-priority preemptive scheduler} with
round-robin inside each priority level. The revised presentation follows the
usual paper pattern of making a claim in prose, then pointing to a numbered
listing that carries the supporting code \citep{symbian_eka2_kernel}.

\paragraph{Scheduling Architecture}

Runnable threads are organised by absolute priority rather than by a fairness
score. EKA2 exposes 64 absolute scheduler priorities in
\texttt{include/nkern/nkern.h}, and Listing~\ref{lst:sym-tprilistbase} shows
that ready threads are stored in per-priority queues backed by a presence
bitmap. That design makes lookup cheap and deterministic, which is exactly the
trade-off Symbian wants for mobile workloads that care more about bounded
latency than about global fairness \citep{sales2005, symbian_eka2_kernel}.

\begin{listing}[H]
\begin{minted}{cpp}
class TPriListBase
	{
public:
	IMPORT_C TPriListBase(TInt aNumPriorities);
	IMPORT_C TInt HighestPriority();
	IMPORT_C TPriListLink* First();
	IMPORT_C void Add(TPriListLink* aLink);
	IMPORT_C void AddHead(TPriListLink* aLink);
	IMPORT_C void Remove(TPriListLink* aLink);
	IMPORT_C void ChangePriority(TPriListLink* aLink, TInt aNewPriority);

	inline TBool NonEmpty() const
		{ return iPresent[0]|iPresent[1]; }
	inline TBool IsEmpty() const
		{ return !NonEmpty();  }
	inline TBool operator>(TInt p) const
		{ return ((p<32) ? (iPresent[1] | (iPresent[0]>>p)>>1)
		                 : (iPresent[1]>>(p-32))>>1 ); }
public:
	union
		{
		TUint iPresent[2];
		TUint64 iPresent64;
		};
	SDblQueLink* iQueue[1];
	};
\end{minted}
\caption{Priority bitmap and ready-queue anchors in
\texttt{include/nkern/nklib.h}.}
\label{lst:sym-tprilistbase}
\end{listing}

The concrete scheduler object reinforces that priority-first design.
Listing~\ref{lst:sym-tscheduler} shows that \texttt{TScheduler} tracks the
current thread, deferred function call queues, delayed work, and reschedule
state rather than a large body of fairness metadata. That contrasts sharply
with Linux's Completely Fair Scheduler (CFS), which is structured around
fairness accounting.

\begin{listing}[H]
\begin{minted}{cpp}
class TScheduler : public TPriListBase
	{
public:
	TScheduler();
	void Remove(NThreadBase* aThread);
	void RotateReadyList(TInt aPriority);
	void QueueDfcs();
	static void Reschedule();
	static void YieldTo(NThreadBase* aThread);
	void TimesliceTick();
	IMPORT_C static TScheduler* Ptr();
	inline void SetProcessHandler(TLinAddr aHandler)
		{iProcessHandler=aHandler;}
private:
	SDblQueLink* iExtraQueues[KNumPriorities-1];
public:
	TUint8 iRescheduleNeededFlag;
	TUint8 iDfcPendingFlag;
	TInt iKernCSLocked;
	SDblQue iDfcs;
	TLinAddr iMonitorExceptionHandler;
	TLinAddr iProcessHandler;
	TLinAddr iRescheduleHook;
	TUint8 iInIDFC;
	NFastMutex iLock;
	NThreadBase* iCurrentThread;
	TAny* iAddressSpace;
	SDblQue iIdleDfcs;
	SDblQue iDelayedQ;
	TDfc iDelayDfc;
	TUint iMadeReadyCounter;
	TUint iMadeUnReadyCounter;
	TUint iTimeSliceExpireCounter;
	};
\end{minted}
\caption{Core scheduler state in \texttt{include/nkern/nk\_priv.h}.}
\label{lst:sym-tscheduler}
\end{listing}

\paragraph{Priority and Timeslice Model}

Symbian does not treat thread priority as one flat application-controlled
number. Instead, it combines an 8-level process-priority band with a relative
thread priority, then maps that pair to an absolute scheduler priority through
the lookup logic in Listing~\ref{lst:sym-priority-table}
\citep{symbian_eka2_kernel}. The table makes the policy explicit: priorities
are clustered into bands, real-time ranges are capability-gated, and a
thread's effective priority depends partly on the role of its owning process.

\begin{listing}[H]
\begin{minted}{cpp}
enum TProcPriority
	{
	EProcPriorityLow=0,
	EProcPriorityBackground=1,
	EProcPriorityForeground=2,
	EProcPriorityHigh=3,
	EProcPrioritySystemServer1=4,
	EProcPrioritySystemServer2=5,
	EProcPrioritySystemServer3=6,
	EProcPriorityRealTimeServer=7,
	};

LOCAL_D const TUint8 ThreadPriorityTable[64] =
	{
/*Low*/               1,     1,     2,     3,  4,  5, 22,  0,
/*Background*/        3,     5,     6,     7,  8,  9, 22,  0,
/*Foreground*/        3,    10,    11,    12, 13, 14, 22,  0,
/*High*/              3,    17,    18,    19, 20, 22, 23,  0,
/*SystemServer1*/     9,    15,    16,    21,
                       KPrioritySystemServerMore, 25, 28, 0,
/*SystemServer2*/     9,    15,    16,    21,
                       KPrioritySystemServerMore, 25, 28, 0,
/*SystemServer3*/     9,    15,    16,    21,
                       KPrioritySystemServerMore, 25, 28, 0,
/*RealTimeServer*/   18,    26,    27,    28, 29, 30, 31, 0
	};

TInt DThread::CalcDefaultThreadPriority()
	{
	TInt r;
	TInt tp=iThreadPriority;
	if (tp>=0)
		r=(tp<KNumPriorities)?tp:KNumPriorities-1;
	else
		{
		tp+=8;
		if (tp<0)
			tp=0;
		TInt pp=iOwningProcess->iPriority;
		TInt i=(pp<<3)+tp;
		r=ThreadPriorityTable[i];
		if ((r > KMaxPriorityWithoutProtServ) &&
			!(this->HasCapability(ECapabilityProtServ,
			  __PLATSEC_DIAGNOSTIC_STRING(
			    "Warning: thread priority capped to below realtime range"))))
			r = KMaxPriorityWithoutProtServ;
		}
	return r;
	}
\end{minted}
\caption{Process-priority bands and thread-priority mapping. Source:
\texttt{include/kernel/kern\_priv.h} and \texttt{kernel/sthread.cpp}.}
\label{lst:sym-priority-table}
\end{listing}

Timeslice handling is equally direct. Listing~\ref{lst:sym-timeslice} combines
the per-thread time accounting fields, the default 20~ms user timeslice, and
the scheduler routines that decrement quanta, restore them on block, and rotate
the ready list among equal-priority peers. The resulting behavior is exactly
what the implementation suggests: CPU-bound peers share by round-robin,
input/output (I/O)-bound threads tend to come back with a fresh quantum, and
lower-priority work can starve if higher-priority threads never block.

\begin{listing}[H]
\begin{minted}{cpp}
TInt iTime;
TInt iTimeslice;

enum {EDefaultUserTimeSliceMs = 20};

void TScheduler::TimesliceTick()
	{
	NThreadBase* pC=iCurrentThread;
	if (pC->iTime>0 && --pC->iTime==0)
		RescheduleNeeded();
	}

void TScheduler::Remove(NThreadBase* aThread)
	{
	__NK_ASSERT_DEBUG(!aThread->iHeldFastMutex);
	iMadeUnReadyCounter++;
	aThread->iTime=aThread->iTimeslice;
	TPriListBase::Remove(aThread);
	}

void TScheduler::RotateReadyList(TInt p)
	{
	__NK_ASSERT_DEBUG(p>=0 && p<KNumPriorities);
	SDblQueLink* pQ=iQueue[p];
	if (pQ)
		{
		SDblQueLink* pN=pQ->iNext;
		if (pN!=pQ)
			{
			NThread* pT=(NThread*)pQ;
			pT->iTime=pT->iTimeslice;
			iQueue[p]=pN;
			if (pQ==iCurrentThread)
				RescheduleNeeded();
			}
		}
	}
\end{minted}
\caption{Timeslice state and round-robin rotation. Source:
\texttt{include/nkern/nk\_priv.h}, \texttt{include/kernel/kern\_priv.h}, and
\texttt{nkern/sched.cpp}.}
\label{lst:sym-timeslice}
\end{listing}

\paragraph{Process and Thread Model}

Symbian uses a \textbf{1:1 kernel-thread model}: every \texttt{DThread}
corresponds to a real nanokernel thread, with no user-level scheduler sitting
above it \citep{symbian_eka2_kernel}. Process-level priority still matters,
however, because changing the process band triggers a recalculation of the
owned threads' effective priorities. Listing~\ref{lst:sym-process-setpriority}
shows that \texttt{DProcess::SetPriority()} walks the process thread list and
recomputes each thread's effective priority.

\begin{listing}[H]
\begin{minted}{cpp}
TInt DProcess::SetPriority(TProcessPriority aPriority)
	{
	TInt p=ConvertPriority(aPriority);
	if (p<0)
		return KErrArgument;
	if (iExitType!=EExitPending)
		return KErrDied;
	__KTRACE_OPT(KPROC,Kern::Printf("Process %O SetPriority(%d)",this,p));
	TInt r=WaitProcessLock();
	__KTRACE_FAIL(r,Kern::Printf("PSP: %d",r));
	if (r!=KErrNone)
		return KErrDied;
#ifdef BTRACE_THREAD_PRIORITY
	BTrace8(BTrace::EThreadPriority,BTrace::EProcessPriority,this,p);
#endif
	NKern::LockSystem();
	iPriority=p;
	SDblQueLink* pLink=iThreadQ.First();
	while (pLink!=&iThreadQ.iA)
		{
		DThread* pT=_LOFF(pLink,DThread,iProcessLink);
		pT->SetThreadPriority(pT->iThreadPriority);
		pLink=pLink->iNext;
		NKern::FlashSystem();
		}
	NKern::UnlockSystem();
	SignalProcessLock();
	return KErrNone;
	}
\end{minted}
\caption{Process-priority propagation in \texttt{kernel/sprocess.cpp}.}
\label{lst:sym-process-setpriority}
\end{listing}

Thread creation is split between the executive and the nanokernel.
Listings~\ref{lst:sym-thread-create-exec} and
\ref{lst:sym-thread-create-nkern} show the handoff clearly: the executive
allocates the supervisor stack and fills out an \texttt{SNThreadCreateInfo},
then the ARM nanokernel constructs the initial saved register frame. The main
advantage is clarity. A blocking kernel operation blocks the real thread
directly, but each thread also consumes actual kernel resources, so the model
scales less gracefully than designs intended for very large thread counts.

\begin{listing}[H]
\begin{minted}{cpp}
TInt DThread::DoCreate(SThreadCreateInfo& aInfo)
	{
	iSupervisorStackSize=aInfo.iSupervisorStackSize;
	if (iSupervisorStackSize==0)
		iSupervisorStackSize=K::SupervisorThreadStackSize;
	if (iThreadType==EThreadInitial || iThreadType==EThreadAPInitial
	 || iThreadType==EThreadMinimalSupervisor)
		iSupervisorStack=aInfo.iSupervisorStack;
	TInt r=KErrNone;
	if (!iSupervisorStack)
		r=AllocateSupervisorStack();
	if (r!=KErrNone)
		return r;
	SNThreadCreateInfo ni;
	ni.iFunction=EpocThreadFunction;
	ni.iPriority=1;
	ni.iTimeslice = (iThreadType==EThreadMinimalSupervisor
	              || iThreadType==EThreadAPInitial
	              || iThreadType==EThreadInitial)
	               ? -1
	               : NKern::TimesliceTicks(EDefaultUserTimeSliceMs*1000);
	ni.iHandlers = &EpocThreadHandlers;
	ni.iFastExecTable=(const SFastExecTable*)EpocFastExecTable;
	ni.iSlowExecTable=(const SSlowExecTable*)EpocSlowExecTable;
	ni.iParameterBlock=(const TUint32*)&aInfo;
	ni.iParameterBlockSize=aInfo.iTotalSize;
	ni.iStackBase=iSupervisorStack;
	ni.iStackSize=iSupervisorStackSize;
	}
\end{minted}
\caption{Executive-side thread setup. Source: \texttt{kernel/sthread.cpp}.}
\label{lst:sym-thread-create-exec}
\end{listing}

\begin{listing}[H]
\begin{minted}{cpp}
TInt NThread::Create(SNThreadCreateInfo& aInfo, TBool aInitial)
	{
	TInt r=NThreadBase::Create(aInfo,aInitial);
	if (r!=KErrNone)
		return r;
	if (!aInitial)
		{
		TUint32* sp=
			(TUint32*)(iStackBase+iStackSize-aInfo.iParameterBlockSize);
		TUint32 r6=(TUint32)aInfo.iParameterBlock;
		if (aInfo.iParameterBlockSize)
			{
			wordmove(sp,aInfo.iParameterBlock,aInfo.iParameterBlockSize);
			r6=(TUint32)sp;
			}
		*--sp=(TUint32)__StartThread;
		*--sp=0;
		*--sp=0;
		*--sp=0;
		*--sp=0;
		*--sp=0;
		*--sp=r6;
		*--sp=(TUint32)aInfo.iFunction;
		*--sp=(TUint32)this;
		*--sp=0x13;
		*--sp=0;
		*--sp=0;
		}
	}
\end{minted}
\caption{Nanokernel thread bootstrap frame construction. Source:
\texttt{nkern/arm/ncthrd.cpp}.}
\label{lst:sym-thread-create-nkern}
\end{listing}

\paragraph{Context Switching}

The context-switch path is tightly coupled to the machine.
Listing~\ref{lst:sym-reschedule} captures the hot reschedule path in
\texttt{nkern/arm/ncsched.cia}: interrupts are masked, pending DFC work is
checked, the current context is saved, and the bitmap-backed ready queues are
consulted to find the highest-priority runnable thread. This is where Symbian's
design priorities are most obvious. The kernel is not trying to build a large
adaptive scheduler; it is trying to get from an interrupt or blocking event to
the right next thread with as little overhead as possible.

\begin{listing}[H]
\begin{minted}{cpp}
__NAKED__ void TScheduler::Reschedule()
	{
	asm("ldr r0, __TheScheduler ");
	asm("str lr, [sp, #-4]! ");
	SET_INTS(r3, MODE_SVC, INTS_ALL_OFF);
	asm("ldrb r1, [r0, #%a0]" : : "i"
		_FOFF(TScheduler,iDfcPendingFlag));
	asm("mov r2, #0 ");
	asm("cmp r1, #0 ");
	asm("start_resched: ");
	asm("blne  " CSM_ZN10TScheduler9QueueDfcsEv);
	asm("ldrb r1, [r0, #%a0]" : : "i"
		_FOFF(TScheduler,iRescheduleNeededFlag));
	SET_INTS_1(r3, MODE_SVC, INTS_ALL_ON);
	asm("cmp r1, #0 ");
	asm("beq no_resched_needed ");
	SET_INTS_2(r3, MODE_SVC, INTS_ALL_ON);
	asm("mrs r2, spsr ");
	asm("ldr r1, [r0, #%a0]" : : "i"
		_FOFF(TScheduler,iCurrentThread));
	asm("stmfd sp!, {r2,r4-r11} ");
	asm("str sp, [r1, #%a0]" : : "i"
		_FOFF(NThread,iSavedSP));
#ifdef __CPU_ARM_HAS_CLZ
	CLZ(12,14);
	asm("subs r12, r12, #32 ");
	CLZcc(CC_EQ,12,1);
	asm("rsb r12, r12, #31 ");
#endif
	asm("ldr r2, [r0, r12, lsl #2] ");
	}
\end{minted}
\caption{Hot reschedule path in \texttt{nkern/arm/ncsched.cia}.}
\label{lst:sym-reschedule}
\end{listing}

\paragraph{Behavior Under Different Workloads}

The workload behavior follows directly from Listings~\ref{lst:sym-priority-table}
and \ref{lst:sym-timeslice}. Equal-priority CPU-bound threads share the
processor by round-robin. I/O-bound threads often perform well because blocking
restores their timeslice. Lower-priority work may still starve if a
higher-priority runnable thread never blocks. Symbian treats that outcome as an
acceptable consequence of explicit priority policy rather than as a scheduling
bug \citep{sales2005, symbian_eka2_kernel}.

\paragraph{Multicore (Symmetric Multiprocessing, SMP) Support}

When built with \texttt{\_\_SMP\_\_}, EKA2 preserves the fixed-priority model
but changes the scheduler topology. Listing~\ref{lst:sym-smp} groups the three
main pieces: a priority-to-class mapping for load balancing, explicit
load-balancing thresholds, and thread-creation logic that records CPU-affinity
or SMP-unsafe compatibility constraints. Symbian therefore adds balancing on
top of strict priorities rather than replacing strict priorities with a fully
fair scheduler.

\begin{listing}[H]
\begin{minted}{cpp}
const TUint8 KClassFromPriority[KNumPriorities] =
	{
	0,0,0,0,0,0,0,0,
	0,0,0,0,1,1,1,1,
	2,2,2,2,2,2,2,2,
	2,2,2,3,3,3,3,3,
	3,3,3,3,3,3,3,3,
	3,3,3,3,3,3,3,3,
	3,3,3,3,3,3,3,3,
	3,3,3,3,3,3,3,3
	};

const TInt  K_LB_BalanceInterval      = 107;
const TInt  K_LB_CpuLoadDiffThreshold = 128;
const TUint K_LB_HeavyPriorityThreshold = 25;

#ifdef __SMP__
	TUint32 config = TheSuperPage().KernelConfigFlags();
	ni.iGroup = 0;
	ni.iCpuAffinity = KCpuAffinityAny;
	if (iThreadType==EThreadUser)
		{
		if ((config & EKernelConfigSMPUnsafeCPU0)
		 && iOwningProcess->iSMPUnsafeCount)
			{
			ni.iCpuAffinity = 0;
			}
		else
			{
			if ((config & EKernelConfigSMPUnsafeCompat)
			 && iOwningProcess->iSMPUnsafeCount)
				ni.iGroup = iOwningProcess->iSMPUnsafeGroup;
			}
		}
#endif
\end{minted}
\caption{SMP priority classes, balance thresholds, and affinity handling.
Source: \texttt{nkernsmp/sched.cpp}, \texttt{nkernsmp/nk\_bal.cpp}, and
\texttt{kernel/sthread.cpp}.}
\label{lst:sym-smp}
\end{listing}

\paragraph{Starvation, Priority Inheritance, and Fairness}

The fixed-priority model inherently permits starvation across priority levels.
To reduce the more damaging problem of priority inversion, EKA2 implements
bounded priority inheritance through the thread cleanup queue.
Listing~\ref{lst:sym-priority-inheritance} shows the key mechanism:
\texttt{SetRequiredPriority()} raises the effective thread priority to the
maximum of the thread's default priority and the highest inherited cleanup
priority. This is a narrow but sensible mitigation. Symbian does not try to
equalise long-term progress across all runnable work; instead, it fixes the
specific failure mode most likely to break a fixed-priority real-time system.

\begin{listing}[H]
\begin{minted}{cpp}
EXPORT_C TThreadCleanup::TThreadCleanup()
	: iThread(NULL)
	{
	}

void TThreadCleanup::ChangePriority(TInt aNewPriority)
	{
	if (aNewPriority != iPriority)
		{
		if (iThread)
			{
			iThread->iCleanupQ.ChangePriority(this,aNewPriority);
			if (++K::PINestLevel<KMaxPriorityInheritanceNesting)
				iThread->SetRequiredPriority();
			}
		else
			iPriority=TUint8(aNewPriority);
		}
	}

EXPORT_C void TThreadCleanup::Remove()
	{
	iThread->iCleanupQ.Remove(this);
	if (iPriority>=iThread->iNThread.i_NThread_BasePri)
		iThread->SetRequiredPriority();
	}

void DThread::SetRequiredPriority()
	{
	TInt p=iDefaultPriority;
	TInt c=iCleanupQ.HighestPriority();
	if (c>p)
		p=c;
	if (p!=iNThread.i_NThread_BasePri)
		SetActualPriority(p);
	}
\end{minted}
\caption{Priority inheritance through cleanup-linked propagation. Source:
\texttt{kernel/sthread.cpp}.}
\label{lst:sym-priority-inheritance}
\end{listing}
